\documentclass[11pt]{article}

    \usepackage[breakable]{tcolorbox}
    \usepackage{parskip} % Stop auto-indenting (to mimic markdown behaviour)
    

    % Basic figure setup, for now with no caption control since it's done
    % automatically by Pandoc (which extracts ![](path) syntax from Markdown).
    \usepackage{graphicx}
    % Keep aspect ratio if custom image width or height is specified
    \setkeys{Gin}{keepaspectratio}
    % Maintain compatibility with old templates. Remove in nbconvert 6.0
    \let\Oldincludegraphics\includegraphics
    % Ensure that by default, figures have no caption (until we provide a
    % proper Figure object with a Caption API and a way to capture that
    % in the conversion process - todo).
    \usepackage{caption}
    \DeclareCaptionFormat{nocaption}{}
    \captionsetup{format=nocaption,aboveskip=0pt,belowskip=0pt}

    \usepackage{float}
    \floatplacement{figure}{H} % forces figures to be placed at the correct location
    \usepackage{xcolor} % Allow colors to be defined
    \usepackage{enumerate} % Needed for markdown enumerations to work
    \usepackage{geometry} % Used to adjust the document margins
    \usepackage{amsmath} % Equations
    \usepackage{amssymb} % Equations
    \usepackage{textcomp} % defines textquotesingle
    % Hack from http://tex.stackexchange.com/a/47451/13684:
    \AtBeginDocument{%
        \def\PYZsq{\textquotesingle}% Upright quotes in Pygmentized code
    }
    \usepackage{upquote} % Upright quotes for verbatim code
    \usepackage{eurosym} % defines \euro

    \usepackage{iftex}
    \ifPDFTeX
        \usepackage[T1]{fontenc}
        \IfFileExists{alphabeta.sty}{
              \usepackage{alphabeta}
          }{
              \usepackage[mathletters]{ucs}
              \usepackage[utf8x]{inputenc}
          }
    \else
        \usepackage{fontspec}
        \usepackage{unicode-math}
    \fi

    \usepackage{fancyvrb} % verbatim replacement that allows latex
    \usepackage{grffile} % extends the file name processing of package graphics
                         % to support a larger range
    \makeatletter % fix for old versions of grffile with XeLaTeX
    \@ifpackagelater{grffile}{2019/11/01}
    {
      % Do nothing on new versions
    }
    {
      \def\Gread@@xetex#1{%
        \IfFileExists{"\Gin@base".bb}%
        {\Gread@eps{\Gin@base.bb}}%
        {\Gread@@xetex@aux#1}%
      }
    }
    \makeatother
    \usepackage[Export]{adjustbox} % Used to constrain images to a maximum size
    \adjustboxset{max size={0.9\linewidth}{0.9\paperheight}}

    % The hyperref package gives us a pdf with properly built
    % internal navigation ('pdf bookmarks' for the table of contents,
    % internal cross-reference links, web links for URLs, etc.)
    \usepackage{hyperref}
    % The default LaTeX title has an obnoxious amount of whitespace. By default,
    % titling removes some of it. It also provides customization options.
    \usepackage{titling}
    \usepackage{longtable} % longtable support required by pandoc >1.10
    \usepackage{booktabs}  % table support for pandoc > 1.12.2
    \usepackage{array}     % table support for pandoc >= 2.11.3
    \usepackage{calc}      % table minipage width calculation for pandoc >= 2.11.1
    \usepackage[inline]{enumitem} % IRkernel/repr support (it uses the enumerate* environment)
    \usepackage[normalem]{ulem} % ulem is needed to support strikethroughs (\sout)
                                % normalem makes italics be italics, not underlines
    \usepackage{soul}      % strikethrough (\st) support for pandoc >= 3.0.0
    \usepackage{mathrsfs}
    

    
    % Colors for the hyperref package
    \definecolor{urlcolor}{rgb}{0,.145,.698}
    \definecolor{linkcolor}{rgb}{.71,0.21,0.01}
    \definecolor{citecolor}{rgb}{.12,.54,.11}

    % ANSI colors
    \definecolor{ansi-black}{HTML}{3E424D}
    \definecolor{ansi-black-intense}{HTML}{282C36}
    \definecolor{ansi-red}{HTML}{E75C58}
    \definecolor{ansi-red-intense}{HTML}{B22B31}
    \definecolor{ansi-green}{HTML}{00A250}
    \definecolor{ansi-green-intense}{HTML}{007427}
    \definecolor{ansi-yellow}{HTML}{DDB62B}
    \definecolor{ansi-yellow-intense}{HTML}{B27D12}
    \definecolor{ansi-blue}{HTML}{208FFB}
    \definecolor{ansi-blue-intense}{HTML}{0065CA}
    \definecolor{ansi-magenta}{HTML}{D160C4}
    \definecolor{ansi-magenta-intense}{HTML}{A03196}
    \definecolor{ansi-cyan}{HTML}{60C6C8}
    \definecolor{ansi-cyan-intense}{HTML}{258F8F}
    \definecolor{ansi-white}{HTML}{C5C1B4}
    \definecolor{ansi-white-intense}{HTML}{A1A6B2}
    \definecolor{ansi-default-inverse-fg}{HTML}{FFFFFF}
    \definecolor{ansi-default-inverse-bg}{HTML}{000000}

    % common color for the border for error outputs.
    \definecolor{outerrorbackground}{HTML}{FFDFDF}

    % commands and environments needed by pandoc snippets
    % extracted from the output of `pandoc -s`
    \providecommand{\tightlist}{%
      \setlength{\itemsep}{0pt}\setlength{\parskip}{0pt}}
    \DefineVerbatimEnvironment{Highlighting}{Verbatim}{commandchars=\\\{\}}
    % Add ',fontsize=\small' for more characters per line
    \newenvironment{Shaded}{}{}
    \newcommand{\KeywordTok}[1]{\textcolor[rgb]{0.00,0.44,0.13}{\textbf{{#1}}}}
    \newcommand{\DataTypeTok}[1]{\textcolor[rgb]{0.56,0.13,0.00}{{#1}}}
    \newcommand{\DecValTok}[1]{\textcolor[rgb]{0.25,0.63,0.44}{{#1}}}
    \newcommand{\BaseNTok}[1]{\textcolor[rgb]{0.25,0.63,0.44}{{#1}}}
    \newcommand{\FloatTok}[1]{\textcolor[rgb]{0.25,0.63,0.44}{{#1}}}
    \newcommand{\CharTok}[1]{\textcolor[rgb]{0.25,0.44,0.63}{{#1}}}
    \newcommand{\StringTok}[1]{\textcolor[rgb]{0.25,0.44,0.63}{{#1}}}
    \newcommand{\CommentTok}[1]{\textcolor[rgb]{0.38,0.63,0.69}{\textit{{#1}}}}
    \newcommand{\OtherTok}[1]{\textcolor[rgb]{0.00,0.44,0.13}{{#1}}}
    \newcommand{\AlertTok}[1]{\textcolor[rgb]{1.00,0.00,0.00}{\textbf{{#1}}}}
    \newcommand{\FunctionTok}[1]{\textcolor[rgb]{0.02,0.16,0.49}{{#1}}}
    \newcommand{\RegionMarkerTok}[1]{{#1}}
    \newcommand{\ErrorTok}[1]{\textcolor[rgb]{1.00,0.00,0.00}{\textbf{{#1}}}}
    \newcommand{\NormalTok}[1]{{#1}}

    % Additional commands for more recent versions of Pandoc
    \newcommand{\ConstantTok}[1]{\textcolor[rgb]{0.53,0.00,0.00}{{#1}}}
    \newcommand{\SpecialCharTok}[1]{\textcolor[rgb]{0.25,0.44,0.63}{{#1}}}
    \newcommand{\VerbatimStringTok}[1]{\textcolor[rgb]{0.25,0.44,0.63}{{#1}}}
    \newcommand{\SpecialStringTok}[1]{\textcolor[rgb]{0.73,0.40,0.53}{{#1}}}
    \newcommand{\ImportTok}[1]{{#1}}
    \newcommand{\DocumentationTok}[1]{\textcolor[rgb]{0.73,0.13,0.13}{\textit{{#1}}}}
    \newcommand{\AnnotationTok}[1]{\textcolor[rgb]{0.38,0.63,0.69}{\textbf{\textit{{#1}}}}}
    \newcommand{\CommentVarTok}[1]{\textcolor[rgb]{0.38,0.63,0.69}{\textbf{\textit{{#1}}}}}
    \newcommand{\VariableTok}[1]{\textcolor[rgb]{0.10,0.09,0.49}{{#1}}}
    \newcommand{\ControlFlowTok}[1]{\textcolor[rgb]{0.00,0.44,0.13}{\textbf{{#1}}}}
    \newcommand{\OperatorTok}[1]{\textcolor[rgb]{0.40,0.40,0.40}{{#1}}}
    \newcommand{\BuiltInTok}[1]{{#1}}
    \newcommand{\ExtensionTok}[1]{{#1}}
    \newcommand{\PreprocessorTok}[1]{\textcolor[rgb]{0.74,0.48,0.00}{{#1}}}
    \newcommand{\AttributeTok}[1]{\textcolor[rgb]{0.49,0.56,0.16}{{#1}}}
    \newcommand{\InformationTok}[1]{\textcolor[rgb]{0.38,0.63,0.69}{\textbf{\textit{{#1}}}}}
    \newcommand{\WarningTok}[1]{\textcolor[rgb]{0.38,0.63,0.69}{\textbf{\textit{{#1}}}}}
    \makeatletter
    \newsavebox\pandoc@box
    \newcommand*\pandocbounded[1]{%
      \sbox\pandoc@box{#1}%
      % scaling factors for width and height
      \Gscale@div\@tempa\textheight{\dimexpr\ht\pandoc@box+\dp\pandoc@box\relax}%
      \Gscale@div\@tempb\linewidth{\wd\pandoc@box}%
      % select the smaller of both
      \ifdim\@tempb\p@<\@tempa\p@
        \let\@tempa\@tempb
      \fi
      % scaling accordingly (\@tempa < 1)
      \ifdim\@tempa\p@<\p@
        \scalebox{\@tempa}{\usebox\pandoc@box}%
      % scaling not needed, use as it is
      \else
        \usebox{\pandoc@box}%
      \fi
    }
    \makeatother

    % Define a nice break command that doesn't care if a line doesn't already
    % exist.
    \def\br{\hspace*{\fill} \\* }
    % Math Jax compatibility definitions
    \def\gt{>}
    \def\lt{<}
    \let\Oldtex\TeX
    \let\Oldlatex\LaTeX
    \renewcommand{\TeX}{\textrm{\Oldtex}}
    \renewcommand{\LaTeX}{\textrm{\Oldlatex}}
    % Document parameters
    % Document title
    \title{Notebook\_to\_build\_graphics}
    
    
    
    
    
    
    
% Pygments definitions
\makeatletter
\def\PY@reset{\let\PY@it=\relax \let\PY@bf=\relax%
    \let\PY@ul=\relax \let\PY@tc=\relax%
    \let\PY@bc=\relax \let\PY@ff=\relax}
\def\PY@tok#1{\csname PY@tok@#1\endcsname}
\def\PY@toks#1+{\ifx\relax#1\empty\else%
    \PY@tok{#1}\expandafter\PY@toks\fi}
\def\PY@do#1{\PY@bc{\PY@tc{\PY@ul{%
    \PY@it{\PY@bf{\PY@ff{#1}}}}}}}
\def\PY#1#2{\PY@reset\PY@toks#1+\relax+\PY@do{#2}}

\@namedef{PY@tok@w}{\def\PY@tc##1{\textcolor[rgb]{0.73,0.73,0.73}{##1}}}
\@namedef{PY@tok@c}{\let\PY@it=\textit\def\PY@tc##1{\textcolor[rgb]{0.24,0.48,0.48}{##1}}}
\@namedef{PY@tok@cp}{\def\PY@tc##1{\textcolor[rgb]{0.61,0.40,0.00}{##1}}}
\@namedef{PY@tok@k}{\let\PY@bf=\textbf\def\PY@tc##1{\textcolor[rgb]{0.00,0.50,0.00}{##1}}}
\@namedef{PY@tok@kp}{\def\PY@tc##1{\textcolor[rgb]{0.00,0.50,0.00}{##1}}}
\@namedef{PY@tok@kt}{\def\PY@tc##1{\textcolor[rgb]{0.69,0.00,0.25}{##1}}}
\@namedef{PY@tok@o}{\def\PY@tc##1{\textcolor[rgb]{0.40,0.40,0.40}{##1}}}
\@namedef{PY@tok@ow}{\let\PY@bf=\textbf\def\PY@tc##1{\textcolor[rgb]{0.67,0.13,1.00}{##1}}}
\@namedef{PY@tok@nb}{\def\PY@tc##1{\textcolor[rgb]{0.00,0.50,0.00}{##1}}}
\@namedef{PY@tok@nf}{\def\PY@tc##1{\textcolor[rgb]{0.00,0.00,1.00}{##1}}}
\@namedef{PY@tok@nc}{\let\PY@bf=\textbf\def\PY@tc##1{\textcolor[rgb]{0.00,0.00,1.00}{##1}}}
\@namedef{PY@tok@nn}{\let\PY@bf=\textbf\def\PY@tc##1{\textcolor[rgb]{0.00,0.00,1.00}{##1}}}
\@namedef{PY@tok@ne}{\let\PY@bf=\textbf\def\PY@tc##1{\textcolor[rgb]{0.80,0.25,0.22}{##1}}}
\@namedef{PY@tok@nv}{\def\PY@tc##1{\textcolor[rgb]{0.10,0.09,0.49}{##1}}}
\@namedef{PY@tok@no}{\def\PY@tc##1{\textcolor[rgb]{0.53,0.00,0.00}{##1}}}
\@namedef{PY@tok@nl}{\def\PY@tc##1{\textcolor[rgb]{0.46,0.46,0.00}{##1}}}
\@namedef{PY@tok@ni}{\let\PY@bf=\textbf\def\PY@tc##1{\textcolor[rgb]{0.44,0.44,0.44}{##1}}}
\@namedef{PY@tok@na}{\def\PY@tc##1{\textcolor[rgb]{0.41,0.47,0.13}{##1}}}
\@namedef{PY@tok@nt}{\let\PY@bf=\textbf\def\PY@tc##1{\textcolor[rgb]{0.00,0.50,0.00}{##1}}}
\@namedef{PY@tok@nd}{\def\PY@tc##1{\textcolor[rgb]{0.67,0.13,1.00}{##1}}}
\@namedef{PY@tok@s}{\def\PY@tc##1{\textcolor[rgb]{0.73,0.13,0.13}{##1}}}
\@namedef{PY@tok@sd}{\let\PY@it=\textit\def\PY@tc##1{\textcolor[rgb]{0.73,0.13,0.13}{##1}}}
\@namedef{PY@tok@si}{\let\PY@bf=\textbf\def\PY@tc##1{\textcolor[rgb]{0.64,0.35,0.47}{##1}}}
\@namedef{PY@tok@se}{\let\PY@bf=\textbf\def\PY@tc##1{\textcolor[rgb]{0.67,0.36,0.12}{##1}}}
\@namedef{PY@tok@sr}{\def\PY@tc##1{\textcolor[rgb]{0.64,0.35,0.47}{##1}}}
\@namedef{PY@tok@ss}{\def\PY@tc##1{\textcolor[rgb]{0.10,0.09,0.49}{##1}}}
\@namedef{PY@tok@sx}{\def\PY@tc##1{\textcolor[rgb]{0.00,0.50,0.00}{##1}}}
\@namedef{PY@tok@m}{\def\PY@tc##1{\textcolor[rgb]{0.40,0.40,0.40}{##1}}}
\@namedef{PY@tok@gh}{\let\PY@bf=\textbf\def\PY@tc##1{\textcolor[rgb]{0.00,0.00,0.50}{##1}}}
\@namedef{PY@tok@gu}{\let\PY@bf=\textbf\def\PY@tc##1{\textcolor[rgb]{0.50,0.00,0.50}{##1}}}
\@namedef{PY@tok@gd}{\def\PY@tc##1{\textcolor[rgb]{0.63,0.00,0.00}{##1}}}
\@namedef{PY@tok@gi}{\def\PY@tc##1{\textcolor[rgb]{0.00,0.52,0.00}{##1}}}
\@namedef{PY@tok@gr}{\def\PY@tc##1{\textcolor[rgb]{0.89,0.00,0.00}{##1}}}
\@namedef{PY@tok@ge}{\let\PY@it=\textit}
\@namedef{PY@tok@gs}{\let\PY@bf=\textbf}
\@namedef{PY@tok@gp}{\let\PY@bf=\textbf\def\PY@tc##1{\textcolor[rgb]{0.00,0.00,0.50}{##1}}}
\@namedef{PY@tok@go}{\def\PY@tc##1{\textcolor[rgb]{0.44,0.44,0.44}{##1}}}
\@namedef{PY@tok@gt}{\def\PY@tc##1{\textcolor[rgb]{0.00,0.27,0.87}{##1}}}
\@namedef{PY@tok@err}{\def\PY@bc##1{{\setlength{\fboxsep}{\string -\fboxrule}\fcolorbox[rgb]{1.00,0.00,0.00}{1,1,1}{\strut ##1}}}}
\@namedef{PY@tok@kc}{\let\PY@bf=\textbf\def\PY@tc##1{\textcolor[rgb]{0.00,0.50,0.00}{##1}}}
\@namedef{PY@tok@kd}{\let\PY@bf=\textbf\def\PY@tc##1{\textcolor[rgb]{0.00,0.50,0.00}{##1}}}
\@namedef{PY@tok@kn}{\let\PY@bf=\textbf\def\PY@tc##1{\textcolor[rgb]{0.00,0.50,0.00}{##1}}}
\@namedef{PY@tok@kr}{\let\PY@bf=\textbf\def\PY@tc##1{\textcolor[rgb]{0.00,0.50,0.00}{##1}}}
\@namedef{PY@tok@bp}{\def\PY@tc##1{\textcolor[rgb]{0.00,0.50,0.00}{##1}}}
\@namedef{PY@tok@fm}{\def\PY@tc##1{\textcolor[rgb]{0.00,0.00,1.00}{##1}}}
\@namedef{PY@tok@vc}{\def\PY@tc##1{\textcolor[rgb]{0.10,0.09,0.49}{##1}}}
\@namedef{PY@tok@vg}{\def\PY@tc##1{\textcolor[rgb]{0.10,0.09,0.49}{##1}}}
\@namedef{PY@tok@vi}{\def\PY@tc##1{\textcolor[rgb]{0.10,0.09,0.49}{##1}}}
\@namedef{PY@tok@vm}{\def\PY@tc##1{\textcolor[rgb]{0.10,0.09,0.49}{##1}}}
\@namedef{PY@tok@sa}{\def\PY@tc##1{\textcolor[rgb]{0.73,0.13,0.13}{##1}}}
\@namedef{PY@tok@sb}{\def\PY@tc##1{\textcolor[rgb]{0.73,0.13,0.13}{##1}}}
\@namedef{PY@tok@sc}{\def\PY@tc##1{\textcolor[rgb]{0.73,0.13,0.13}{##1}}}
\@namedef{PY@tok@dl}{\def\PY@tc##1{\textcolor[rgb]{0.73,0.13,0.13}{##1}}}
\@namedef{PY@tok@s2}{\def\PY@tc##1{\textcolor[rgb]{0.73,0.13,0.13}{##1}}}
\@namedef{PY@tok@sh}{\def\PY@tc##1{\textcolor[rgb]{0.73,0.13,0.13}{##1}}}
\@namedef{PY@tok@s1}{\def\PY@tc##1{\textcolor[rgb]{0.73,0.13,0.13}{##1}}}
\@namedef{PY@tok@mb}{\def\PY@tc##1{\textcolor[rgb]{0.40,0.40,0.40}{##1}}}
\@namedef{PY@tok@mf}{\def\PY@tc##1{\textcolor[rgb]{0.40,0.40,0.40}{##1}}}
\@namedef{PY@tok@mh}{\def\PY@tc##1{\textcolor[rgb]{0.40,0.40,0.40}{##1}}}
\@namedef{PY@tok@mi}{\def\PY@tc##1{\textcolor[rgb]{0.40,0.40,0.40}{##1}}}
\@namedef{PY@tok@il}{\def\PY@tc##1{\textcolor[rgb]{0.40,0.40,0.40}{##1}}}
\@namedef{PY@tok@mo}{\def\PY@tc##1{\textcolor[rgb]{0.40,0.40,0.40}{##1}}}
\@namedef{PY@tok@ch}{\let\PY@it=\textit\def\PY@tc##1{\textcolor[rgb]{0.24,0.48,0.48}{##1}}}
\@namedef{PY@tok@cm}{\let\PY@it=\textit\def\PY@tc##1{\textcolor[rgb]{0.24,0.48,0.48}{##1}}}
\@namedef{PY@tok@cpf}{\let\PY@it=\textit\def\PY@tc##1{\textcolor[rgb]{0.24,0.48,0.48}{##1}}}
\@namedef{PY@tok@c1}{\let\PY@it=\textit\def\PY@tc##1{\textcolor[rgb]{0.24,0.48,0.48}{##1}}}
\@namedef{PY@tok@cs}{\let\PY@it=\textit\def\PY@tc##1{\textcolor[rgb]{0.24,0.48,0.48}{##1}}}

\def\PYZbs{\char`\\}
\def\PYZus{\char`\_}
\def\PYZob{\char`\{}
\def\PYZcb{\char`\}}
\def\PYZca{\char`\^}
\def\PYZam{\char`\&}
\def\PYZlt{\char`\<}
\def\PYZgt{\char`\>}
\def\PYZsh{\char`\#}
\def\PYZpc{\char`\%}
\def\PYZdl{\char`\$}
\def\PYZhy{\char`\-}
\def\PYZsq{\char`\'}
\def\PYZdq{\char`\"}
\def\PYZti{\char`\~}
% for compatibility with earlier versions
\def\PYZat{@}
\def\PYZlb{[}
\def\PYZrb{]}
\makeatother


    % For linebreaks inside Verbatim environment from package fancyvrb.
    \makeatletter
        \newbox\Wrappedcontinuationbox
        \newbox\Wrappedvisiblespacebox
        \newcommand*\Wrappedvisiblespace {\textcolor{red}{\textvisiblespace}}
        \newcommand*\Wrappedcontinuationsymbol {\textcolor{red}{\llap{\tiny$\m@th\hookrightarrow$}}}
        \newcommand*\Wrappedcontinuationindent {3ex }
        \newcommand*\Wrappedafterbreak {\kern\Wrappedcontinuationindent\copy\Wrappedcontinuationbox}
        % Take advantage of the already applied Pygments mark-up to insert
        % potential linebreaks for TeX processing.
        %        {, <, #, %, $, ' and ": go to next line.
        %        _, }, ^, &, >, - and ~: stay at end of broken line.
        % Use of \textquotesingle for straight quote.
        \newcommand*\Wrappedbreaksatspecials {%
            \def\PYGZus{\discretionary{\char`\_}{\Wrappedafterbreak}{\char`\_}}%
            \def\PYGZob{\discretionary{}{\Wrappedafterbreak\char`\{}{\char`\{}}%
            \def\PYGZcb{\discretionary{\char`\}}{\Wrappedafterbreak}{\char`\}}}%
            \def\PYGZca{\discretionary{\char`\^}{\Wrappedafterbreak}{\char`\^}}%
            \def\PYGZam{\discretionary{\char`\&}{\Wrappedafterbreak}{\char`\&}}%
            \def\PYGZlt{\discretionary{}{\Wrappedafterbreak\char`\<}{\char`\<}}%
            \def\PYGZgt{\discretionary{\char`\>}{\Wrappedafterbreak}{\char`\>}}%
            \def\PYGZsh{\discretionary{}{\Wrappedafterbreak\char`\#}{\char`\#}}%
            \def\PYGZpc{\discretionary{}{\Wrappedafterbreak\char`\%}{\char`\%}}%
            \def\PYGZdl{\discretionary{}{\Wrappedafterbreak\char`\$}{\char`\$}}%
            \def\PYGZhy{\discretionary{\char`\-}{\Wrappedafterbreak}{\char`\-}}%
            \def\PYGZsq{\discretionary{}{\Wrappedafterbreak\textquotesingle}{\textquotesingle}}%
            \def\PYGZdq{\discretionary{}{\Wrappedafterbreak\char`\"}{\char`\"}}%
            \def\PYGZti{\discretionary{\char`\~}{\Wrappedafterbreak}{\char`\~}}%
        }
        % Some characters . , ; ? ! / are not pygmentized.
        % This macro makes them "active" and they will insert potential linebreaks
        \newcommand*\Wrappedbreaksatpunct {%
            \lccode`\~`\.\lowercase{\def~}{\discretionary{\hbox{\char`\.}}{\Wrappedafterbreak}{\hbox{\char`\.}}}%
            \lccode`\~`\,\lowercase{\def~}{\discretionary{\hbox{\char`\,}}{\Wrappedafterbreak}{\hbox{\char`\,}}}%
            \lccode`\~`\;\lowercase{\def~}{\discretionary{\hbox{\char`\;}}{\Wrappedafterbreak}{\hbox{\char`\;}}}%
            \lccode`\~`\:\lowercase{\def~}{\discretionary{\hbox{\char`\:}}{\Wrappedafterbreak}{\hbox{\char`\:}}}%
            \lccode`\~`\?\lowercase{\def~}{\discretionary{\hbox{\char`\?}}{\Wrappedafterbreak}{\hbox{\char`\?}}}%
            \lccode`\~`\!\lowercase{\def~}{\discretionary{\hbox{\char`\!}}{\Wrappedafterbreak}{\hbox{\char`\!}}}%
            \lccode`\~`\/\lowercase{\def~}{\discretionary{\hbox{\char`\/}}{\Wrappedafterbreak}{\hbox{\char`\/}}}%
            \catcode`\.\active
            \catcode`\,\active
            \catcode`\;\active
            \catcode`\:\active
            \catcode`\?\active
            \catcode`\!\active
            \catcode`\/\active
            \lccode`\~`\~
        }
    \makeatother

    \let\OriginalVerbatim=\Verbatim
    \makeatletter
    \renewcommand{\Verbatim}[1][1]{%
        %\parskip\z@skip
        \sbox\Wrappedcontinuationbox {\Wrappedcontinuationsymbol}%
        \sbox\Wrappedvisiblespacebox {\FV@SetupFont\Wrappedvisiblespace}%
        \def\FancyVerbFormatLine ##1{\hsize\linewidth
            \vtop{\raggedright\hyphenpenalty\z@\exhyphenpenalty\z@
                \doublehyphendemerits\z@\finalhyphendemerits\z@
                \strut ##1\strut}%
        }%
        % If the linebreak is at a space, the latter will be displayed as visible
        % space at end of first line, and a continuation symbol starts next line.
        % Stretch/shrink are however usually zero for typewriter font.
        \def\FV@Space {%
            \nobreak\hskip\z@ plus\fontdimen3\font minus\fontdimen4\font
            \discretionary{\copy\Wrappedvisiblespacebox}{\Wrappedafterbreak}
            {\kern\fontdimen2\font}%
        }%

        % Allow breaks at special characters using \PYG... macros.
        \Wrappedbreaksatspecials
        % Breaks at punctuation characters . , ; ? ! and / need catcode=\active
        \OriginalVerbatim[#1,codes*=\Wrappedbreaksatpunct]%
    }
    \makeatother

    % Exact colors from NB
    \definecolor{incolor}{HTML}{303F9F}
    \definecolor{outcolor}{HTML}{D84315}
    \definecolor{cellborder}{HTML}{CFCFCF}
    \definecolor{cellbackground}{HTML}{F7F7F7}

    % prompt
    \makeatletter
    \newcommand{\boxspacing}{\kern\kvtcb@left@rule\kern\kvtcb@boxsep}
    \makeatother
    \newcommand{\prompt}[4]{
        {\ttfamily\llap{{\color{#2}[#3]:\hspace{3pt}#4}}\vspace{-\baselineskip}}
    }
    

    
    % Prevent overflowing lines due to hard-to-break entities
    \sloppy
    % Setup hyperref package
    \hypersetup{
      breaklinks=true,  % so long urls are correctly broken across lines
      colorlinks=true,
      urlcolor=urlcolor,
      linkcolor=linkcolor,
      citecolor=citecolor,
      }
    % Slightly bigger margins than the latex defaults
    
    \geometry{verbose,tmargin=1in,bmargin=1in,lmargin=1in,rmargin=1in}
    
    

\begin{document}
    
    \maketitle
    
    

    
    \hypertarget{introduction}{%
\section{Introduction}\label{introduction}}

    \textbf{File:} Notebook\_to\_build\_graphics.ipynb\\
\textbf{Author:} Elizabeth Gould\\
\textbf{Date:} 10.03.2026\\
\textbf{Problem:} Here there are two: 1) Plot \(\psi(x)\) 2) Test our
model for \(M_{\mu}\left(BR, dk, \mu, \psi_0^\prime\right)\)

    This code contains two short exercises. The first is just a sample of a
run of our derivative grid code, which plots the IVP solutions for
\(\psi(x)\) The second is a short Monte Carlo test of our model for
\(M_{\mu}\left(BR, dk, \mu, \psi_0^\prime\right)\), which contains both
visual and numerical analysis.

    \hypertarget{derivative-grids}{%
\section{Derivative Grids}\label{derivative-grids}}

    This is the code used for building the derivative grid graphs with the
deriv\_grid class of my library.

    \begin{tcolorbox}[breakable, size=fbox, boxrule=1pt, pad at break*=1mm,colback=cellbackground, colframe=cellborder]
\prompt{In}{incolor}{2}{\boxspacing}
\begin{Verbatim}[commandchars=\\\{\}]
\PY{c+c1}{\PYZsh{}\PYZhy{}\PYZhy{}LIBRARIES\PYZhy{}\PYZhy{}\PYZhy{}\PYZhy{}\PYZhy{}\PYZhy{}\PYZhy{}\PYZhy{}}

\PY{c+c1}{\PYZsh{}numerics}
\PY{k+kn}{import} \PY{n+nn}{numpy} \PY{k}{as} \PY{n+nn}{np}
\PY{k+kn}{import} \PY{n+nn}{eelib}

\PY{c+c1}{\PYZsh{}\PYZhy{}\PYZhy{}CODE\PYZhy{}\PYZhy{}}

\PY{c+c1}{\PYZsh{} set parameters}
\PY{c+c1}{\PYZsh{} note that k, B, and R are percents here, mu is not as its scale is unknown}
\PY{n}{k} \PY{o}{=} \PY{l+m+mf}{0.5}
\PY{n}{R} \PY{o}{=} \PY{l+m+mf}{1.0}
\PY{n}{B} \PY{o}{=} \PY{l+m+mf}{0.5}
\PY{n}{mu} \PY{o}{=} \PY{l+m+mf}{1.0e\PYZhy{}7}

\PY{n}{n\PYZus{}g} \PY{o}{=} \PY{l+m+mi}{5}

\PY{n}{pr} \PY{o}{=} \PY{l+m+mi}{1000}
\PY{n}{n} \PY{o}{=} \PY{l+m+mi}{1}

\PY{n}{loopl} \PY{o}{=} \PY{n}{eelib}\PY{o}{.}\PY{n}{deriv\PYZus{}grid}\PY{p}{(}\PY{n}{R}\PY{p}{,} \PY{n}{B}\PY{p}{,} \PY{n}{k}\PY{p}{,} \PY{n}{mu}\PY{p}{,} \PY{n}{grid\PYZus{}size}\PY{o}{=}\PY{n}{n\PYZus{}g}\PY{p}{)}
\PY{n}{loopl}\PY{o}{.}\PY{n}{derivGrid}\PY{p}{(}\PY{p}{)}


\PY{c+c1}{\PYZsh{} \PYZhy{}\PYZhy{}\PYZhy{} PLOTTING FOR GRIDS \PYZhy{}\PYZhy{}\PYZhy{}}
\PY{k}{for} \PY{n}{i} \PY{o+ow}{in} \PY{n+nb}{range}\PY{p}{(}\PY{n}{n\PYZus{}g}\PY{p}{)}\PY{p}{:}
    \PY{k}{for} \PY{n}{j} \PY{o+ow}{in} \PY{n+nb}{range}\PY{p}{(}\PY{n}{n\PYZus{}g}\PY{p}{)}\PY{p}{:}
        \PY{n}{loopl}\PY{o}{.}\PY{n}{plot\PYZus{}abs}\PY{p}{(}\PY{n}{i}\PY{p}{,}\PY{n}{j}\PY{p}{,}\PY{n}{k}\PY{p}{,}\PY{n}{R}\PY{p}{,}\PY{n}{B}\PY{p}{,}\PY{n}{mu}\PY{p}{,}\PY{n}{n}\PY{p}{)}
        \PY{n}{loopl}\PY{o}{.}\PY{n}{plot\PYZus{}real}\PY{p}{(}\PY{n}{i}\PY{p}{,}\PY{n}{j}\PY{p}{,}\PY{n}{k}\PY{p}{,}\PY{n}{R}\PY{p}{,}\PY{n}{B}\PY{p}{,}\PY{n}{mu}\PY{p}{,}\PY{n}{n}\PY{p}{)}
\end{Verbatim}
\end{tcolorbox}

    \begin{Verbatim}[commandchars=\\\{\}]
Begin grid build:  0.0
Plot Code: ['er', '0r', '0x', 'em'] 35
i =  1 of  5 time:  0.0010137557983398438
Done, i, j =  1 1 time:  425.60918831825256
Done, i, j =  1 2 time:  999.7543675899506
Done, i, j =  1 3 time:  1670.3892183303833
Done, i, j =  1 4 time:  2243.107450723648
Done, i, j =  1 5 time:  2865.3794503211975
i =  2 of  5 time:  2865.394219636917
Done, i, j =  2 1 time:  3425.1825251579285
Done, i, j =  2 2 time:  3972.909587621689
Done, i, j =  2 3 time:  4369.62785243988
Done, i, j =  2 4 time:  4774.924133777618
Done, i, j =  2 5 time:  5237.2201454639435
i =  3 of  5 time:  5237.236102104187
Done, i, j =  3 1 time:  5627.000266313553
Done, i, j =  3 2 time:  6041.5580015182495
Done, i, j =  3 3 time:  6338.41552567482
Done, i, j =  3 4 time:  6732.538671731949
Done, i, j =  3 5 time:  7079.623293161392
i =  4 of  5 time:  7079.645211935043
Done, i, j =  4 1 time:  7499.045262813568
Done, i, j =  4 2 time:  7839.223477840424
Done, i, j =  4 3 time:  8120.520606517792
Done, i, j =  4 4 time:  8450.375850915909
Done, i, j =  4 5 time:  8783.337816953659
i =  5 of  5 time:  8783.34081530571
Done, i, j =  5 1 time:  9141.17084312439
Done, i, j =  5 2 time:  9511.636178016663
Done, i, j =  5 3 time:  9992.288011550903
Done, i, j =  5 4 time:  10488.71034526825
Done, i, j =  5 5 time:  10984.962453603745
Done grid build:  10985.002418518066
    \end{Verbatim}

    \hypertarget{test-of-slow-oscillation-wavenumber-model}{%
\section{Test of Slow Oscillation Wavenumber
Model}\label{test-of-slow-oscillation-wavenumber-model}}

    This is code to visually and numerically test the fit of my slow
oscillation model. The model is not perfect, but good enough.

    \begin{tcolorbox}[breakable, size=fbox, boxrule=1pt, pad at break*=1mm,colback=cellbackground, colframe=cellborder]
\prompt{In}{incolor}{ }{\boxspacing}
\begin{Verbatim}[commandchars=\\\{\}]
\PY{k+kn}{import} \PY{n+nn}{eelib}
\PY{k+kn}{import} \PY{n+nn}{numpy} \PY{k}{as} \PY{n+nn}{np}
\PY{k+kn}{import} \PY{n+nn}{matplotlib}\PY{n+nn}{.}\PY{n+nn}{pyplot} \PY{k}{as} \PY{n+nn}{plt}
\end{Verbatim}
\end{tcolorbox}

    \begin{tcolorbox}[breakable, size=fbox, boxrule=1pt, pad at break*=1mm,colback=cellbackground, colframe=cellborder]
\prompt{In}{incolor}{2}{\boxspacing}
\begin{Verbatim}[commandchars=\\\{\}]
\PY{c+c1}{\PYZsh{}\PYZhy{}\PYZhy{}PARAMETERS\PYZhy{}\PYZhy{}\PYZhy{}\PYZhy{}\PYZhy{}\PYZhy{}\PYZhy{}\PYZhy{}}

\PY{c+c1}{\PYZsh{} note that k, B, and R are percents here, mu is not as its scale is unknown}
\PY{n}{dk} \PY{o}{=} \PY{l+m+mf}{0.5}
\PY{n}{R}  \PY{o}{=} \PY{l+m+mf}{1.0}
\PY{n}{B}  \PY{o}{=} \PY{l+m+mf}{0.5}
\PY{n}{mu} \PY{o}{=} \PY{l+m+mf}{1.0e\PYZhy{}6}

\PY{n}{n\PYZus{}mc} \PY{o}{=} \PY{l+m+mi}{10}

\PY{n}{b\PYZus{}r}  \PY{o}{=} \PY{p}{(}\PY{l+m+mf}{0.5}\PY{p}{,} \PY{l+m+mf}{1.0}\PY{p}{)}  \PY{c+c1}{\PYZsh{} B here is B*R}
\PY{n}{mu\PYZus{}r} \PY{o}{=} \PY{p}{(}\PY{o}{\PYZhy{}}\PY{l+m+mi}{8}\PY{p}{,} \PY{o}{\PYZhy{}}\PY{l+m+mi}{6}\PY{p}{)}    \PY{c+c1}{\PYZsh{} 1.0e\PYZhy{}8 to 1.0e\PYZhy{}6}
\end{Verbatim}
\end{tcolorbox}

    \begin{tcolorbox}[breakable, size=fbox, boxrule=1pt, pad at break*=1mm,colback=cellbackground, colframe=cellborder]
\prompt{In}{incolor}{3}{\boxspacing}
\begin{Verbatim}[commandchars=\\\{\}]
\PY{c+c1}{\PYZsh{}\PYZhy{}\PYZhy{}CODE\PYZhy{}\PYZhy{}}

\PY{c+c1}{\PYZsh{} Make the grid object}
\PY{n}{gridl} \PY{o}{=} \PY{n}{eelib}\PY{o}{.}\PY{n}{grid\PYZus{}slow\PYZus{}osc}\PY{p}{(}\PY{n}{R}\PY{p}{,} \PY{n}{B}\PY{p}{,} \PY{n}{dk}\PY{p}{,} \PY{n}{mu}\PY{p}{)}

\PY{c+c1}{\PYZsh{} Save our full solution}
\PY{n}{gridl}\PY{o}{.}\PY{n}{save\PYZus{}solution} \PY{o}{=} \PY{k+kc}{True}

\PY{c+c1}{\PYZsh{} Make the grid}
\PY{n}{gridl}\PY{o}{.}\PY{n}{makeMCPoints}\PY{p}{(}\PY{n}{mu}\PY{o}{=}\PY{n}{mu\PYZus{}r}\PY{p}{,} \PY{n}{B}\PY{o}{=}\PY{n}{b\PYZus{}r}\PY{p}{,} \PY{n}{num} \PY{o}{=} \PY{n}{n\PYZus{}mc}\PY{p}{)}

\PY{c+c1}{\PYZsh{} Run the grid}
\PY{n}{gridl}\PY{o}{.}\PY{n}{mcSlowOsc}\PY{p}{(}\PY{p}{)}
\end{Verbatim}
\end{tcolorbox}

    \begin{Verbatim}[commandchars=\\\{\}]
Begin grid build:  0.0
Number of periods to calculate: 10
Done grid build:  1436.7411334514618
    \end{Verbatim}

    \begin{tcolorbox}[breakable, size=fbox, boxrule=1pt, pad at break*=1mm,colback=cellbackground, colframe=cellborder]
\prompt{In}{incolor}{25}{\boxspacing}
\begin{Verbatim}[commandchars=\\\{\}]
\PY{c+c1}{\PYZsh{} First, declare which solution I am using}
\PY{n}{ii}  \PY{o}{=} \PY{n+nb}{int}\PY{p}{(}\PY{n}{np}\PY{o}{.}\PY{n}{random}\PY{o}{.}\PY{n}{rand}\PY{p}{(}\PY{p}{)}\PY{o}{*}\PY{n}{n\PYZus{}mc}\PY{p}{)}
\PY{k}{if} \PY{n}{ii} \PY{o}{\PYZgt{}}\PY{o}{=} \PY{n}{n\PYZus{}mc}\PY{p}{:} \PY{n}{ii} \PY{o}{=} \PY{n}{n\PYZus{}mc} \PY{o}{\PYZhy{}} \PY{l+m+mi}{1} \PY{c+c1}{\PYZsh{} just in case}
\PY{n}{sol} \PY{o}{=} \PY{n}{gridl}\PY{o}{.}\PY{n}{slow\PYZus{}osc\PYZus{}sol}\PY{p}{[}\PY{n}{ii}\PY{p}{]}

\PY{c+c1}{\PYZsh{} Retrieve my fits for this function from the grid.}
\PY{n}{MM}    \PY{o}{=} \PY{n}{gridl}\PY{o}{.}\PY{n}{slow\PYZus{}osc\PYZus{}k}\PY{p}{[}\PY{n}{ii}\PY{p}{]}
\PY{n}{amp}   \PY{o}{=} \PY{n}{gridl}\PY{o}{.}\PY{n}{slow\PYZus{}osc\PYZus{}a}\PY{p}{[}\PY{n}{ii}\PY{p}{]}
\PY{n}{theta} \PY{o}{=} \PY{n}{gridl}\PY{o}{.}\PY{n}{slow\PYZus{}osc\PYZus{}th}\PY{p}{[}\PY{n}{ii}\PY{p}{]}

\PY{c+c1}{\PYZsh{} Predicted slow oscillation amplitude.}
\PY{n}{vt} \PY{o}{=} \PY{n}{gridl}\PY{o}{.}\PY{n}{val\PYZus{}table}\PY{p}{[}\PY{n}{ii}\PY{p}{]} \PY{c+c1}{\PYZsh{} mu, dk, B, R, A, k0, dr, di}
\PY{n}{MM2} \PY{o}{=} \PY{n}{eelib}\PY{o}{.}\PY{n}{pred\PYZus{}slow\PYZus{}k}\PY{p}{(}\PY{n}{vt}\PY{p}{[}\PY{l+m+mi}{6}\PY{p}{]}\PY{o}{+}\PY{l+m+mi}{1}\PY{n}{j}\PY{o}{*}\PY{n}{vt}\PY{p}{[}\PY{l+m+mi}{7}\PY{p}{]}\PY{p}{,} \PY{n}{vt}\PY{p}{[}\PY{l+m+mi}{0}\PY{p}{]}\PY{p}{,} \PY{n}{vt}\PY{p}{[}\PY{l+m+mi}{0}\PY{p}{]}\PY{p}{,} \PY{n}{vt}\PY{p}{[}\PY{l+m+mi}{2}\PY{p}{]}\PY{o}{*}\PY{n}{eelib}\PY{o}{.}\PY{n}{B\PYZus{}max}\PY{p}{,} \PY{n}{vt}\PY{p}{[}\PY{l+m+mi}{3}\PY{p}{]}\PY{o}{*}\PY{n}{eelib}\PY{o}{.}\PY{n}{R\PYZus{}max}\PY{p}{,} \PY{n}{vt}\PY{p}{[}\PY{l+m+mi}{4}\PY{p}{]}\PY{p}{,} \PY{n}{vt}\PY{p}{[}\PY{l+m+mi}{5}\PY{p}{]}\PY{p}{)} \PY{c+c1}{\PYZsh{}dpsi0, mu, dk, B, R, A = 1., k0=kFAu}

\PY{c+c1}{\PYZsh{} Pull the found solution from the sol variable.}
\PY{n}{t\PYZus{}list}  \PY{o}{=} \PY{n}{sol}\PY{p}{[}\PY{l+s+s1}{\PYZsq{}}\PY{l+s+s1}{t}\PY{l+s+s1}{\PYZsq{}}\PY{p}{]}
\PY{n}{y1\PYZus{}list} \PY{o}{=} \PY{n}{np}\PY{o}{.}\PY{n}{real}\PY{p}{(}\PY{n}{sol}\PY{p}{[}\PY{l+s+s1}{\PYZsq{}}\PY{l+s+s1}{y}\PY{l+s+s1}{\PYZsq{}}\PY{p}{]}\PY{p}{[}\PY{l+m+mi}{0}\PY{p}{]}\PY{p}{)}

\PY{c+c1}{\PYZsh{} Estimated fit}
\PY{n}{y2\PYZus{}list} \PY{o}{=} \PY{n}{amp} \PY{o}{*} \PY{n}{np}\PY{o}{.}\PY{n}{sin}\PY{p}{(}\PY{n}{MM2} \PY{o}{*} \PY{p}{(}\PY{n}{t\PYZus{}list}\PY{p}{)} \PY{o}{+} \PY{n}{theta}\PY{p}{)}

\PY{c+c1}{\PYZsh{} Fit to a sin using SciPy}
\PY{n}{y3\PYZus{}list} \PY{o}{=} \PY{n}{amp} \PY{o}{*} \PY{n}{np}\PY{o}{.}\PY{n}{sin}\PY{p}{(}\PY{n}{MM} \PY{o}{*} \PY{n}{t\PYZus{}list} \PY{o}{+} \PY{n}{theta}\PY{p}{)}
\end{Verbatim}
\end{tcolorbox}

    \begin{tcolorbox}[breakable, size=fbox, boxrule=1pt, pad at break*=1mm,colback=cellbackground, colframe=cellborder]
\prompt{In}{incolor}{ }{\boxspacing}
\begin{Verbatim}[commandchars=\\\{\}]
\PY{c+c1}{\PYZsh{} And plot the results}
\PY{n}{fig}\PY{p}{,} \PY{n}{ax} \PY{o}{=} \PY{n}{plt}\PY{o}{.}\PY{n}{subplots}\PY{p}{(}\PY{p}{)}

\PY{n}{ax}\PY{o}{.}\PY{n}{set\PYZus{}ylabel}\PY{p}{(}\PY{l+s+s1}{\PYZsq{}}\PY{l+s+s1}{Re(ψ)}\PY{l+s+s1}{\PYZsq{}}\PY{p}{)}
\PY{n}{ax}\PY{o}{.}\PY{n}{set\PYZus{}xlabel}\PY{p}{(}\PY{l+s+s1}{\PYZsq{}}\PY{l+s+s1}{x (m)}\PY{l+s+s1}{\PYZsq{}}\PY{p}{)}
\PY{n}{plt}\PY{o}{.}\PY{n}{title}\PY{p}{(}\PY{l+s+sa}{f}\PY{l+s+s2}{\PYZdq{}}\PY{l+s+s2}{Re(ψ(x)), following maxima of |ψ(x)|}\PY{l+s+s2}{\PYZdq{}}\PY{p}{)}

\PY{n}{line1}\PY{p}{,} \PY{o}{=} \PY{n}{ax}\PY{o}{.}\PY{n}{plot}\PY{p}{(}\PY{n}{t\PYZus{}list}\PY{p}{,} \PY{n}{y1\PYZus{}list}\PY{p}{,} \PY{n}{color} \PY{o}{=} \PY{l+s+s1}{\PYZsq{}}\PY{l+s+s1}{red}\PY{l+s+s1}{\PYZsq{}}\PY{p}{,} \PY{n}{label} \PY{o}{=} \PY{l+s+s1}{\PYZsq{}}\PY{l+s+s1}{Numerical Solution}\PY{l+s+s1}{\PYZsq{}}\PY{p}{)}
\PY{n}{line2}\PY{p}{,} \PY{o}{=} \PY{n}{ax}\PY{o}{.}\PY{n}{plot}\PY{p}{(}\PY{n}{t\PYZus{}list}\PY{p}{,} \PY{n}{y2\PYZus{}list}\PY{p}{,} \PY{n}{color} \PY{o}{=} \PY{l+s+s1}{\PYZsq{}}\PY{l+s+s1}{blue}\PY{l+s+s1}{\PYZsq{}}\PY{p}{,} \PY{n}{label} \PY{o}{=} \PY{l+s+s1}{\PYZsq{}}\PY{l+s+s1}{Predicted Fit}\PY{l+s+s1}{\PYZsq{}}\PY{p}{)}
\PY{n}{line3}\PY{p}{,} \PY{o}{=} \PY{n}{ax}\PY{o}{.}\PY{n}{plot}\PY{p}{(}\PY{n}{t\PYZus{}list}\PY{p}{,} \PY{n}{y3\PYZus{}list}\PY{p}{,} \PY{n}{color} \PY{o}{=} \PY{l+s+s1}{\PYZsq{}}\PY{l+s+s1}{green}\PY{l+s+s1}{\PYZsq{}}\PY{p}{,} \PY{n}{label} \PY{o}{=} \PY{l+s+s1}{\PYZsq{}}\PY{l+s+s1}{Nonlinear Regression Fit}\PY{l+s+s1}{\PYZsq{}}\PY{p}{)}

\PY{n}{ax}\PY{o}{.}\PY{n}{legend}\PY{p}{(}\PY{n}{handles}\PY{o}{=}\PY{p}{[}\PY{n}{line1}\PY{p}{,}\PY{n}{line2}\PY{p}{,}\PY{n}{line3}\PY{p}{]}\PY{p}{)}
\PY{n}{ax}\PY{o}{.}\PY{n}{set\PYZus{}box\PYZus{}aspect}\PY{p}{(}\PY{l+m+mf}{2.0}\PY{o}{/}\PY{l+m+mf}{3.5}\PY{p}{)}

\PY{n}{plt}\PY{o}{.}\PY{n}{show}\PY{p}{(}\PY{p}{)}
\end{Verbatim}
\end{tcolorbox}

    \begin{center}
    \adjustimage{max size={0.9\linewidth}{0.9\paperheight}}{figs/slow_k_ivp_en/output_test_slow_1.png}
    \end{center}
    { \hspace*{\fill} \\}
    
    \begin{tcolorbox}[breakable, size=fbox, boxrule=1pt, pad at break*=1mm,colback=cellbackground, colframe=cellborder]
\prompt{In}{incolor}{ }{\boxspacing}
\begin{Verbatim}[commandchars=\\\{\}]
\PY{c+c1}{\PYZsh{} Plot differences}
\PY{n}{fig}\PY{p}{,} \PY{n}{ax} \PY{o}{=} \PY{n}{plt}\PY{o}{.}\PY{n}{subplots}\PY{p}{(}\PY{p}{)}

\PY{n}{ax}\PY{o}{.}\PY{n}{set\PYZus{}ylabel}\PY{p}{(}\PY{l+s+s1}{\PYZsq{}}\PY{l+s+s1}{ψ(x) Difference}\PY{l+s+s1}{\PYZsq{}}\PY{p}{)}
\PY{n}{ax}\PY{o}{.}\PY{n}{set\PYZus{}xlabel}\PY{p}{(}\PY{l+s+s1}{\PYZsq{}}\PY{l+s+s1}{x (m)}\PY{l+s+s1}{\PYZsq{}}\PY{p}{)}
\PY{n}{plt}\PY{o}{.}\PY{n}{title}\PY{p}{(}\PY{l+s+sa}{f}\PY{l+s+s2}{\PYZdq{}}\PY{l+s+s2}{Differences of the calculated Re(ψ(x)), following maxima of |ψ(x)|}\PY{l+s+s2}{\PYZdq{}}\PY{p}{)}

\PY{n}{line1}\PY{p}{,} \PY{o}{=} \PY{n}{ax}\PY{o}{.}\PY{n}{plot}\PY{p}{(}\PY{n}{t\PYZus{}list}\PY{p}{,} \PY{n}{y1\PYZus{}list}\PY{o}{\PYZhy{}}\PY{n}{y2\PYZus{}list}\PY{p}{,} \PY{n}{color} \PY{o}{=} \PY{l+s+s1}{\PYZsq{}}\PY{l+s+s1}{red}\PY{l+s+s1}{\PYZsq{}}\PY{p}{,} \PY{n}{label} \PY{o}{=} \PY{l+s+s1}{\PYZsq{}}\PY{l+s+s1}{Difference between numerical solution and predicted fit}\PY{l+s+s1}{\PYZsq{}}\PY{p}{)}
\PY{n}{line2}\PY{p}{,} \PY{o}{=} \PY{n}{ax}\PY{o}{.}\PY{n}{plot}\PY{p}{(}\PY{n}{t\PYZus{}list}\PY{p}{,} \PY{n}{y2\PYZus{}list}\PY{o}{\PYZhy{}}\PY{n}{y3\PYZus{}list}\PY{p}{,} \PY{n}{color} \PY{o}{=} \PY{l+s+s1}{\PYZsq{}}\PY{l+s+s1}{blue}\PY{l+s+s1}{\PYZsq{}}\PY{p}{,} \PY{n}{label} \PY{o}{=} \PY{l+s+s1}{\PYZsq{}}\PY{l+s+s1}{Difference between predicted fit and regression fit}\PY{l+s+s1}{\PYZsq{}}\PY{p}{)}
\PY{n}{line3}\PY{p}{,} \PY{o}{=} \PY{n}{ax}\PY{o}{.}\PY{n}{plot}\PY{p}{(}\PY{n}{t\PYZus{}list}\PY{p}{,} \PY{n}{y3\PYZus{}list}\PY{o}{\PYZhy{}}\PY{n}{y1\PYZus{}list}\PY{p}{,} \PY{n}{color} \PY{o}{=} \PY{l+s+s1}{\PYZsq{}}\PY{l+s+s1}{green}\PY{l+s+s1}{\PYZsq{}}\PY{p}{,} \PY{n}{label} \PY{o}{=} \PY{l+s+s1}{\PYZsq{}}\PY{l+s+s1}{Difference between regression fit and numerical solution}\PY{l+s+s1}{\PYZsq{}}\PY{p}{)}

\PY{n}{ax}\PY{o}{.}\PY{n}{legend}\PY{p}{(}\PY{n}{handles}\PY{o}{=}\PY{p}{[}\PY{n}{line1}\PY{p}{,}\PY{n}{line2}\PY{p}{,}\PY{n}{line3}\PY{p}{]}\PY{p}{)}
\PY{n}{ax}\PY{o}{.}\PY{n}{set\PYZus{}box\PYZus{}aspect}\PY{p}{(}\PY{l+m+mf}{2.0}\PY{o}{/}\PY{l+m+mf}{3.5}\PY{p}{)}

\PY{n}{plt}\PY{o}{.}\PY{n}{show}\PY{p}{(}\PY{p}{)}
\end{Verbatim}
\end{tcolorbox}

    \begin{center}
    \adjustimage{max size={0.9\linewidth}{0.9\paperheight}}{figs/slow_k_ivp_en/output_test_slow_2.png}
    \end{center}
    { \hspace*{\fill} \\}
    
    \begin{tcolorbox}[breakable, size=fbox, boxrule=1pt, pad at break*=1mm,colback=cellbackground, colframe=cellborder]
\prompt{In}{incolor}{ }{\boxspacing}
\begin{Verbatim}[commandchars=\\\{\}]
\PY{c+c1}{\PYZsh{} Print the differences}
\PY{n+nb}{print}\PY{p}{(}\PY{l+s+sa}{f}\PY{l+s+s2}{\PYZdq{}}\PY{l+s+s2}{Estimated slow oscillation wave number: }\PY{l+s+si}{\PYZob{}}\PY{n}{MM2}\PY{l+s+si}{\PYZcb{}}\PY{l+s+s2}{\PYZdq{}}\PY{p}{)}
\PY{n+nb}{print}\PY{p}{(}\PY{l+s+sa}{f}\PY{l+s+s2}{\PYZdq{}}\PY{l+s+s2}{Fit slow oscillation wave number: }\PY{l+s+si}{\PYZob{}}\PY{n}{MM}\PY{l+s+si}{\PYZcb{}}\PY{l+s+s2}{\PYZdq{}}\PY{p}{)}
\PY{n+nb}{print}\PY{p}{(}\PY{l+s+s2}{\PYZdq{}}\PY{l+s+s2}{Maximum difference between regression fit and predicted fit:}\PY{l+s+s2}{\PYZdq{}}\PY{p}{,} \PY{n}{np}\PY{o}{.}\PY{n}{max}\PY{p}{(}\PY{n}{y3\PYZus{}list}\PY{o}{\PYZhy{}}\PY{n}{y2\PYZus{}list}\PY{p}{)}\PY{p}{)}
\PY{n+nb}{print}\PY{p}{(}\PY{l+s+s2}{\PYZdq{}}\PY{l+s+s2}{Maximum difference between numerical solution and predicted fit:}\PY{l+s+s2}{\PYZdq{}}\PY{p}{,} \PY{n}{np}\PY{o}{.}\PY{n}{max}\PY{p}{(}\PY{n}{y1\PYZus{}list}\PY{o}{\PYZhy{}}\PY{n}{y2\PYZus{}list}\PY{p}{)}\PY{p}{)}
\PY{n+nb}{print}\PY{p}{(}\PY{l+s+s2}{\PYZdq{}}\PY{l+s+s2}{Maximum difference between numerical solution and regression fit:}\PY{l+s+s2}{\PYZdq{}}\PY{p}{,} \PY{n}{np}\PY{o}{.}\PY{n}{max}\PY{p}{(}\PY{n}{y1\PYZus{}list}\PY{o}{\PYZhy{}}\PY{n}{y3\PYZus{}list}\PY{p}{)}\PY{p}{)}
\end{Verbatim}
\end{tcolorbox}

    \begin{Verbatim}[commandchars=\\\{\}]
Estimated slow oscillation wavenumber: 8389300.433092475
Fit slow oscillation wavenumber: 8387761.299876148
Maximum difference between regression fit and predicted fit:
0.013657253183702513
Maximum difference between numerical solution and predicted fit:
0.01481728630082968
Maximum difference between numerical solution and regression fit:
0.0034084309286432912
    \end{Verbatim}

    \begin{tcolorbox}[breakable, size=fbox, boxrule=1pt, pad at break*=1mm,colback=cellbackground, colframe=cellborder]
\prompt{In}{incolor}{ }{\boxspacing}
\begin{Verbatim}[commandchars=\\\{\}]
\PY{c+c1}{\PYZsh{} Now for the rest of the numbers}
\PY{k}{for} \PY{n}{ii} \PY{o+ow}{in} \PY{n+nb}{range}\PY{p}{(}\PY{n}{n\PYZus{}mc}\PY{p}{)}\PY{p}{:}
    \PY{c+c1}{\PYZsh{} First, declare which solution I am using}
    \PY{n}{sol} \PY{o}{=} \PY{n}{gridl}\PY{o}{.}\PY{n}{slow\PYZus{}osc\PYZus{}sol}\PY{p}{[}\PY{n}{ii}\PY{p}{]}

    \PY{c+c1}{\PYZsh{} Retrieve my fits for this function from the grid.}
    \PY{n}{MM}    \PY{o}{=} \PY{n}{gridl}\PY{o}{.}\PY{n}{slow\PYZus{}osc\PYZus{}k}\PY{p}{[}\PY{n}{ii}\PY{p}{]}
    \PY{n}{amp}   \PY{o}{=} \PY{n}{gridl}\PY{o}{.}\PY{n}{slow\PYZus{}osc\PYZus{}a}\PY{p}{[}\PY{n}{ii}\PY{p}{]}
    \PY{n}{theta} \PY{o}{=} \PY{n}{gridl}\PY{o}{.}\PY{n}{slow\PYZus{}osc\PYZus{}th}\PY{p}{[}\PY{n}{ii}\PY{p}{]}

    \PY{c+c1}{\PYZsh{} Predicted slow oscillation amplitude.}
    \PY{n}{vt} \PY{o}{=} \PY{n}{gridl}\PY{o}{.}\PY{n}{val\PYZus{}table}\PY{p}{[}\PY{n}{ii}\PY{p}{]} \PY{c+c1}{\PYZsh{} mu, dk, B, R, A, k0, dr, di}
    \PY{n}{MM2} \PY{o}{=} \PY{n}{eelib}\PY{o}{.}\PY{n}{pred\PYZus{}slow\PYZus{}k}\PY{p}{(}\PY{n}{vt}\PY{p}{[}\PY{l+m+mi}{6}\PY{p}{]}\PY{o}{+}\PY{l+m+mi}{1}\PY{n}{j}\PY{o}{*}\PY{n}{vt}\PY{p}{[}\PY{l+m+mi}{7}\PY{p}{]}\PY{p}{,} \PY{n}{vt}\PY{p}{[}\PY{l+m+mi}{0}\PY{p}{]}\PY{p}{,} \PY{n}{vt}\PY{p}{[}\PY{l+m+mi}{0}\PY{p}{]}\PY{p}{,} \PY{n}{vt}\PY{p}{[}\PY{l+m+mi}{2}\PY{p}{]}\PY{o}{*}\PY{n}{eelib}\PY{o}{.}\PY{n}{B\PYZus{}max}\PY{p}{,} \PY{n}{vt}\PY{p}{[}\PY{l+m+mi}{3}\PY{p}{]}\PY{o}{*}\PY{n}{eelib}\PY{o}{.}\PY{n}{R\PYZus{}max}\PY{p}{,} \PY{n}{vt}\PY{p}{[}\PY{l+m+mi}{4}\PY{p}{]}\PY{p}{,} \PY{n}{vt}\PY{p}{[}\PY{l+m+mi}{5}\PY{p}{]}\PY{p}{)} \PY{c+c1}{\PYZsh{}dpsi0, mu, dk, B, R, A = 1., k0=kFAu}

    \PY{c+c1}{\PYZsh{} Pull the found solution from the sol variable.}
    \PY{n}{t\PYZus{}list}  \PY{o}{=} \PY{n}{sol}\PY{p}{[}\PY{l+s+s1}{\PYZsq{}}\PY{l+s+s1}{t}\PY{l+s+s1}{\PYZsq{}}\PY{p}{]}
    \PY{n}{y1\PYZus{}list} \PY{o}{=} \PY{n}{np}\PY{o}{.}\PY{n}{real}\PY{p}{(}\PY{n}{sol}\PY{p}{[}\PY{l+s+s1}{\PYZsq{}}\PY{l+s+s1}{y}\PY{l+s+s1}{\PYZsq{}}\PY{p}{]}\PY{p}{[}\PY{l+m+mi}{0}\PY{p}{]}\PY{p}{)}

    \PY{c+c1}{\PYZsh{} Estimated fit}
    \PY{n}{y2\PYZus{}list} \PY{o}{=} \PY{n}{amp} \PY{o}{*} \PY{n}{np}\PY{o}{.}\PY{n}{sin}\PY{p}{(}\PY{n}{MM2} \PY{o}{*} \PY{p}{(}\PY{n}{t\PYZus{}list}\PY{p}{)} \PY{o}{+} \PY{n}{theta}\PY{p}{)}

    \PY{c+c1}{\PYZsh{} Fit to a sin using SciPy}
    \PY{n}{y3\PYZus{}list} \PY{o}{=} \PY{n}{amp} \PY{o}{*} \PY{n}{np}\PY{o}{.}\PY{n}{sin}\PY{p}{(}\PY{n}{MM} \PY{o}{*} \PY{n}{t\PYZus{}list} \PY{o}{+} \PY{n}{theta}\PY{p}{)}

    \PY{c+c1}{\PYZsh{} Print the differences}
    \PY{n+nb}{print}\PY{p}{(}\PY{l+s+sa}{f}\PY{l+s+s2}{\PYZdq{}}\PY{l+s+s2}{Random solution }\PY{l+s+si}{\PYZob{}}\PY{n}{ii}\PY{l+s+si}{\PYZcb{}}\PY{l+s+s2}{\PYZdq{}}\PY{p}{)}
    \PY{n+nb}{print}\PY{p}{(}\PY{l+s+sa}{f}\PY{l+s+s2}{\PYZdq{}}\PY{l+s+s2}{Estimated slow oscillation wave number: }\PY{l+s+si}{\PYZob{}}\PY{n}{MM2}\PY{l+s+si}{\PYZcb{}}\PY{l+s+s2}{\PYZdq{}}\PY{p}{)}
    \PY{n+nb}{print}\PY{p}{(}\PY{l+s+sa}{f}\PY{l+s+s2}{\PYZdq{}}\PY{l+s+s2}{Fit slow oscillation wave number: }\PY{l+s+si}{\PYZob{}}\PY{n}{MM}\PY{l+s+si}{\PYZcb{}}\PY{l+s+s2}{\PYZdq{}}\PY{p}{)}
    \PY{n+nb}{print}\PY{p}{(}\PY{l+s+s2}{\PYZdq{}}\PY{l+s+s2}{Maximum difference between regression fit and estimated fit:}\PY{l+s+s2}{\PYZdq{}}\PY{p}{,} \PY{n}{np}\PY{o}{.}\PY{n}{max}\PY{p}{(}\PY{n}{y3\PYZus{}list}\PY{o}{\PYZhy{}}\PY{n}{y2\PYZus{}list}\PY{p}{)}\PY{p}{)}
    \PY{n+nb}{print}\PY{p}{(}\PY{l+s+s2}{\PYZdq{}}\PY{l+s+s2}{Maximum difference between numerical solution and estimated fit:}\PY{l+s+s2}{\PYZdq{}}\PY{p}{,} \PY{n}{np}\PY{o}{.}\PY{n}{max}\PY{p}{(}\PY{n}{y1\PYZus{}list}\PY{o}{\PYZhy{}}\PY{n}{y2\PYZus{}list}\PY{p}{)}\PY{p}{)}
    \PY{n+nb}{print}\PY{p}{(}\PY{l+s+s2}{\PYZdq{}}\PY{l+s+s2}{Maximum difference between numerical solution and regression fit:}\PY{l+s+s2}{\PYZdq{}}\PY{p}{,} \PY{n}{np}\PY{o}{.}\PY{n}{max}\PY{p}{(}\PY{n}{y1\PYZus{}list}\PY{o}{\PYZhy{}}\PY{n}{y3\PYZus{}list}\PY{p}{)}\PY{p}{)}
    \PY{n+nb}{print}\PY{p}{(}\PY{l+s+s2}{\PYZdq{}}\PY{l+s+s2}{\PYZdq{}}\PY{p}{)}
\end{Verbatim}
\end{tcolorbox}

    \begin{Verbatim}[commandchars=\\\{\}]
Random solution 0
Estimated slow oscillation wavenumber: 7695931.0389159145
Fit slow oscillation wavenumber: 7705523.901747796
Maximum difference between regression fit and estimated fit: 0.08855178629225899
Maximum difference between numerical solution and estimated fit:
0.09008886877059123
Maximum difference between numerical solution and regression fit:
0.003557607193861667

Random solution 1
Estimated slow oscillation wavenumber: 5520152.083362719
Fit slow oscillation wavenumber: 5521561.8786500525
Maximum difference between regression fit and estimated fit:
0.007871258831689008
Maximum difference between numerical solution and estimated fit:
0.007910053928197436
Maximum difference between numerical solution and regression fit:
0.0011458076251886773

Random solution 2
Estimated slow oscillation wavenumber: 10306056.140563946
Fit slow oscillation wavenumber: 10313521.234371202
Maximum difference between regression fit and estimated fit: 0.07582130778922336
Maximum difference between numerical solution and estimated fit:
0.07686257952445752
Maximum difference between numerical solution and regression fit:
0.004369446302817215

Random solution 3
Estimated slow oscillation wavenumber: 5957411.965353246
Fit slow oscillation wavenumber: 5958404.644749804
Maximum difference between regression fit and estimated fit:
0.006004884369374379
Maximum difference between numerical solution and estimated fit:
0.006065564545724447
Maximum difference between numerical solution and regression fit:
0.0014954619400019409

Random solution 4
Estimated slow oscillation wavenumber: 7241263.402560295
Fit slow oscillation wavenumber: 7251959.208602868
Maximum difference between regression fit and estimated fit: 0.08892861931645601
Maximum difference between numerical solution and estimated fit:
0.09004273921871074
Maximum difference between numerical solution and regression fit:
0.0025892565951213253

Random solution 5
Estimated slow oscillation wavenumber: 8502094.345244821
Fit slow oscillation wavenumber: 8517021.835766362
Maximum difference between regression fit and estimated fit: 0.1348280853315711
Maximum difference between numerical solution and estimated fit:
0.13631923711776076
Maximum difference between numerical solution and regression fit:
0.006539412789031029

Random solution 6
Estimated slow oscillation wavenumber: 8249309.549836851
Fit slow oscillation wavenumber: 8242281.861319513
Maximum difference between regression fit and estimated fit: 0.05599140256029245
Maximum difference between numerical solution and estimated fit:
0.056966769712377796
Maximum difference between numerical solution and regression fit:
0.002338612677215135

Random solution 7
Estimated slow oscillation wavenumber: 6868381.901407971
Fit slow oscillation wavenumber: 6873923.747426414
Maximum difference between regression fit and estimated fit: 0.05860081881197013
Maximum difference between numerical solution and estimated fit:
0.058602402146968324
Maximum difference between numerical solution and regression fit:
0.005376837920478161

Random solution 8
Estimated slow oscillation wavenumber: 8389300.433092475
Fit slow oscillation wavenumber: 8387761.299876148
Maximum difference between regression fit and estimated fit:
0.013657253183702513
Maximum difference between numerical solution and estimated fit:
0.01481728630082968
Maximum difference between numerical solution and regression fit:
0.0034084309286432912

Random solution 9
Estimated slow oscillation wavenumber: 8410916.001752483
Fit slow oscillation wavenumber: 8417511.06578477
Maximum difference between regression fit and estimated fit: 0.04041158883247903
Maximum difference between numerical solution and estimated fit:
0.04107973685779891
Maximum difference between numerical solution and regression fit:
0.0013783627551383626

    \end{Verbatim}


    % Add a bibliography block to the postdoc
    
    
    
\end{document}
